\begin{titlepage}
\thispagestyle{empty}
\begin{tikzpicture}[remember picture,overlay]
  \fill[IngBlue] (current page.north west) rectangle ([yshift=-43mm]current page.north east);
  \fill[IngTeal] ([yshift=-43mm]current page.north west) rectangle ([yshift=-47mm]current page.north east);
\end{tikzpicture}
\vspace*{10mm}
{\color{white}\sffamily\bfseries\Large ingegnerismo.it\par}
\vspace{30mm}
{\sffamily\bfseries\fontsize{30}{34}\selectfont\color{IngBlue}\FormularioTitolo\par}
\vspace{5mm}
{\sffamily\fontsize{16}{20}\selectfont\color{IngTeal}\FormularioSottotitolo\par}
\vspace{12mm}
\begin{tcolorbox}[colback=IngPaleBlue,colframe=IngBlue,boxrule=.6pt,arc=1.5mm]
\textbf{Obiettivo.} Individuare rapidamente il tipo di problema, scegliere il metodo adatto, applicare una procedura controllabile e verificare il risultato. Il volume non sostituisce un manuale teorico: raccoglie la teoria strettamente necessaria per usare correttamente formule, criteri e teoremi.
\end{tcolorbox}
\vfill
\begin{tabular}{@{}ll@{}}
\textbf{Autore:} & \FormularioAutore\\
\textbf{Versione:} & \FormularioVersione\\
\textbf{Data:} & \FormularioData\\
\textbf{Sorgente:} & progetto LaTeX versionato\\
\textbf{Sito:} & \href{\SitoURL}{ingegnerismo.it}
\end{tabular}
\vspace{14mm}

{\small\color{IngGray}© 2026 \FormularioAutore. Tutti i diritti riservati. La disponibilità pubblica del repository non attribuisce automaticamente una licenza di riuso.}
\end{titlepage}
